%% 自動生成（gen.py）。直接編集しないこと
\makeatletter
\bkdefvol{1}{かぞえる・はかる}{}{かずで くらべる、かずで つたえる}{小学1年〜3年}{1}{1}
\bkdefvol{2}{\R{数}{かず}と\R{量}{りょう}のことば}{}{「5」ではなく「5\,m」。\R{数}{かず}と\R{量}{りょう}のことばを\R{身}{み}につけよう}{小学4年〜中学1年}{1}{2}
\bkdefvol{3}{\R{関数}{かん|すう}とデータ}{きまりを\R{見}{み}つける}{\R{表}{ひょう}・グラフ・\R{式}{しき}で、\R{変化}{へん|か}のきまりをつかもう}{中学1年〜3年}{1}{3}
\bkdefvol{4}{\R{形}{かたち}と\R{向}{む}き}{\R{図形}{ず|けい}・\R{三角関数}{さん|かく|かん|すう}・ベクトル}{長さと角度で、形と向きを数にしよう}{中学2年〜高校2年}{1}{5}
\bkdefvol{5}{\R{変化}{へん|か}をつかむ}{\R{微分}{び|ぶん}と\R{積分}{せき|ぶん}}{変わる速さと、\mbox{たまった量}を計算しよう}{高校1年〜2年}{2}{4}
\bkdefvol{6}{\R{未来}{み|らい}を\R{計算}{けい|さん}する}{\R{微分方程式}{び|ぶん|ほう|てい|しき}と\R{複素数}{ふく|そ|すう}}{「次の瞬間のルール」から、未来を計算しよう}{高校2年〜3年}{2}{5}
\bkdefvol{7}{まとめて\R{計算}{けい|さん}する}{\R{行列}{ぎょう|れつ}}{たくさんの数を、表のまま計算しよう}{高校3年（発展）〜大学1年}{3}{5}
\bkdefvol{8}{\R{空間}{くう|かん}に\R{広}{ひろ}がる\R{量}{りょう}}{ベクトル\R{解析}{かい|せき}}{空間に広がる量を、矢印で読みとこう}{大学1年〜2年}{3}{5}
\bkdefvol{9}{\R{波}{なみ}と\R{変化}{へん|か}を\R{読}{よ}みとく}{フーリエ\R{解析}{かい|せき}とラプラス\R{変換}{へん|かん}}{波を成分に分け、微分方程式をかけ算で解こう}{大学2年}{3}{5}
\@namedef{bk@rname@1}{\R{変化}{へん|か}}\@namedef{bk@rsub@1}{ひき\R{算}{ざん}→\R{傾}{かたむ}き→\R{微分}{び|ぶん}}\@namedef{bk@rend@1}{\R{流体}{りゅう|たい}・\R{熱}{ねつ}}
\@namedef{bk@rname@2}{たし\R{合}{あ}わせ}\@namedef{bk@rsub@2}{かけ\R{算}{ざん}→\R{面積}{めん|せき}→\R{積分}{せき|ぶん}}\@namedef{bk@rend@2}{\R{電磁気}{でん|じ|き}}
\@namedef{bk@rname@3}{\R{倍々}{ばい|ばい}}\@namedef{bk@rsub@3}{\R{累乗}{るい|じょう}→\R{指数}{し|すう}→$e$}\@namedef{bk@rend@3}{\R{制御}{せい|ぎょ}}
\@namedef{bk@rname@4}{\R{回転}{かい|てん}と\R{波}{なみ}}\@namedef{bk@rsub@4}{\R{角}{かく}→sin→\R{複素数}{ふく|そ|すう}}\@namedef{bk@rend@4}{\R{音}{おと}・\R{交流}{こう|りゅう}}
\@namedef{bk@rname@5}{\R{表}{ひょう}}\@namedef{bk@rsub@5}{\R{九九}{く|く}の\R{表}{ひょう}→\R{行列}{ぎょう|れつ}}\@namedef{bk@rend@5}{\R{回路網}{かい|ろ|もう}・\R{制御}{せい|ぎょ}}
\@namedef{bk@rname@6}{\R{矢印}{や|じるし}}\@namedef{bk@rsub@6}{\R{向}{む}き→ベクトル→\R{場}{ば}}\@namedef{bk@rend@6}{ロボット・\\\R{電磁気}{でん|じ|き}}
\@namedef{bk@rname@7}{はかる}\@namedef{bk@rsub@7}{\R{目}{め}もり→\R{誤差}{ご|さ}→\R{統計}{とう|けい}}\@namedef{bk@rend@7}{\R{計測}{けい|そく}・\R{診断}{しん|だん}}
\makeatother
